\documentclass[aspectratio=169,11pt]{beamer}
% Shared CMPSC 480 Beamer preamble.
\usepackage[T1]{fontenc}
\usepackage[utf8]{inputenc}
\usepackage{lmodern}
\usepackage{amsmath,amssymb,mathtools}
\usepackage{booktabs,array,tabularx}
\usepackage{xcolor}
\usepackage{listings}
\usepackage{tikz}
\usetikzlibrary{arrows.meta,positioning,shapes.geometric}

% Ignore only negligible TeX box diagnostics that are below visual tolerance.
% Larger layout defects still appear in the build log and in rendered QA.
\vfuzz=3pt
\hbadness=2000

\definecolor{courseink}{HTML}{111111}
\definecolor{coursegray}{HTML}{EDEDED}
\definecolor{courserule}{HTML}{B8BCC4}
\definecolor{courseblue}{HTML}{3D8DFF}
\definecolor{courselightblue}{HTML}{D0EDFA}
\definecolor{coursemuted}{HTML}{5C6470}
\definecolor{coursegreen}{HTML}{0B7A53}
\definecolor{coursered}{HTML}{B3261E}

\usetheme{default}
\usefonttheme{professionalfonts}
\setbeamertemplate{navigation symbols}{}
\setbeamersize{text margin left=0.65cm,text margin right=0.65cm}

\setbeamercolor{normal text}{fg=courseink,bg=white}
\setbeamercolor{frametitle}{fg=courseink,bg=white}
\setbeamercolor{title}{fg=courseink,bg=white}
\setbeamercolor{subtitle}{fg=coursemuted,bg=white}
\setbeamercolor{structure}{fg=courseblue}
\setbeamercolor{alerted text}{fg=coursered}
\setbeamercolor{block title}{fg=courseink,bg=courselightblue}
\setbeamercolor{block body}{fg=courseink,bg=coursegray}
\setbeamercolor{example text}{fg=coursegreen}

\setbeamerfont{title}{size=\fontsize{25}{29}\selectfont,series=\bfseries}
\setbeamerfont{subtitle}{size=\fontsize{13}{16}\selectfont}
\setbeamerfont{frametitle}{size=\fontsize{19}{22}\selectfont,series=\bfseries}
\setbeamerfont{normal text}{size=\fontsize{11}{14}\selectfont}
\setbeamerfont{footline}{size=\fontsize{7.5}{9}\selectfont}

\setbeamertemplate{frametitle}{%
  \nointerlineskip
  % Use content-driven height so a long assessment title can wrap without
  % being clipped by a fixed-height title bar.
  \begin{beamercolorbox}[wd=\paperwidth,sep=0.17cm,leftskip=0.48cm,rightskip=0.48cm]{frametitle}
    \insertframetitle
  \end{beamercolorbox}
  \vspace{-0.08cm}
  {\color{courserule}\rule{\textwidth}{0.35pt}}
}

\setbeamertemplate{footline}{%
  \leavevmode
  \begin{beamercolorbox}[wd=\paperwidth,ht=2.6ex,dp=1.1ex,leftskip=.65cm,rightskip=.65cm]{author in head/foot}
    \color{coursemuted}\insertshorttitle\hfill\insertframenumber/\inserttotalframenumber
  \end{beamercolorbox}
}

\setbeamertemplate{itemize item}{\color{courseblue}\small$\blacktriangleright$}
\setbeamertemplate{itemize subitem}{\color{courseblue}\scriptsize$\bullet$}
\setbeamertemplate{enumerate item}{\color{courseblue}\insertenumlabel.}

\lstdefinestyle{coursepython}{
  language=Python,
  basicstyle=\ttfamily\fontsize{8.5}{10}\selectfont,
  keywordstyle=\color{courseblue}\bfseries,
  commentstyle=\color{coursemuted}\itshape,
  stringstyle=\color{coursegreen},
  showstringspaces=false,
  columns=fullflexible,
  keepspaces=true,
  frame=single,
  rulecolor=\color{courserule},
  backgroundcolor=\color{coursegray},
  breaklines=true,
  tabsize=4,
  xleftmargin=0.15cm,
  xrightmargin=0.15cm,
  framexleftmargin=0.15cm,
  framexrightmargin=0.15cm
}
\lstset{style=coursepython}

\newcommand{\points}[1]{\hfill{\color{courseblue}\textbf{[#1 points]}}}
\newcommand{\deliverable}[1]{\textbf{\color{courseblue}#1}}
\newcommand{\code}[1]{\texttt{#1}}
\newcommand{\keyidea}[1]{%
  \begin{beamercolorbox}[rounded=false,sep=0.55em,wd=\linewidth]{block body}
    \textbf{Key idea.} #1
  \end{beamercolorbox}
}
\newcommand{\answerlabel}{\textbf{\color{coursegreen}Answer.}\ }
\newcommand{\gradinglabel}{\textbf{\color{courseblue}Grading guidance.}\ }

\newenvironment{questionblock}[1]
  {\begin{block}{#1}}
  {\end{block}}

\newenvironment{solutionblock}[1]
  {\begin{block}{#1}}
  {\end{block}}

\AtBeginSection[]{
  \begin{frame}
    \vfill
    {\usebeamerfont{title}\color{courseink}\insertsectionhead\par}
    \vspace{0.4cm}
    {\color{courseblue}\rule{0.32\paperwidth}{3pt}}
    \vfill
  \end{frame}
}


% CMPSC480_TOTAL_POINTS: 20
% CMPSC480_ITEM_IDS: 1,2,3,4,5
\title[Module Project A Questions]{Week 3 Module Project A: Search and Adversarial Agents}
\subtitle{Integrating pathfinding with minimax, alpha-beta pruning, and evaluation}
\author{CMPSC 480 - Student Question Deck}
\date{}

\begin{document}


\begin{frame}
  \titlepage
\end{frame}

\begin{frame}{Assessment overview}
\begin{columns}[T,onlytextwidth]
\begin{column}{0.62\textwidth}
\Large Build a deterministic agent that plans paths and chooses actions against an opponent.

\vspace{0.35cm}
\normalsize
Coverage: Weeks 1-3: search formulation, A-star, minimax, alpha-beta, and evaluation
\end{column}
\begin{column}{0.33\textwidth}
\begin{block}{Points}20 total\end{block}
\begin{block}{Expected time}10-14 hours\end{block}
\begin{block}{Submission}Repository, tests, and engineering report\end{block}
\end{column}
\end{columns}
\end{frame}

\begin{frame}{Learning objectives}
\begin{itemize}
  \item Reuse a stable search core for deterministic pathfinding.
  \item Implement depth-limited minimax and alpha-beta pruning.
  \item Design a state evaluation function with explicit feature semantics.
  \item Keep planning and adversarial reasoning behind separate interfaces.
  \item Measure correctness, pruning effectiveness, and decision quality.
\end{itemize}
\end{frame}

\begin{frame}{Requirements checklist}
\small
\begin{itemize}
  \item Use Python 3.11 and NumPy only unless the prompt explicitly permits another package.
  \item Preserve the published interfaces and make all tie-breaking deterministic.
  \item State assumptions, validate inputs, and handle empty, terminal, and unreachable cases.
  \item Include focused tests whose names explain the defect each test would expose.
  \item Report measured evidence; do not replace experimental results with unsupported claims.
  \item Do not mutate game states shared by sibling branches.
  \item Return a legal action for every nonterminal state with at least one legal move.
\end{itemize}
\end{frame}

\begin{frame}{Task 1 - Architecture and pathfinding (4 points)}
\small
Define component boundaries and integrate BFS, DFS, and A-star pathfinding.

\vspace{0.2cm}
\begin{enumerate}
\item[\textbf{a.}] Draw the boundary among game state, search problem, pathfinder, and adversarial agent. \hfill{\color{courseblue}\textbf{[1 points]}}
\item[\textbf{b.}] Implement BFS, DFS, and A-star behind one pathfinding interface. \hfill{\color{courseblue}\textbf{[2 points]}}
\item[\textbf{c.}] Demonstrate path validity and deterministic tie-breaking. \hfill{\color{courseblue}\textbf{[1 points]}}
\end{enumerate}
\end{frame}

\begin{frame}{Task 2 - A-star heuristic (3 points)}
\small
Design, justify, and test a domain-specific heuristic.

\vspace{0.2cm}
\begin{enumerate}
\item[\textbf{a.}] Define the heuristic and its units. \hfill{\color{courseblue}\textbf{[1 points]}}
\item[\textbf{b.}] Give an admissibility or consistency proof sketch. \hfill{\color{courseblue}\textbf{[1 points]}}
\item[\textbf{c.}] Compare A-star with h=0 on at least three maps. \hfill{\color{courseblue}\textbf{[1 points]}}
\end{enumerate}
\end{frame}

\begin{frame}{Task 3 - Minimax decision engine (5 points)}
\small
Implement deterministic, depth-limited minimax with correct terminal semantics.

\vspace{0.2cm}
\begin{enumerate}
\item[\textbf{a.}] Write minimax pseudocode with MAX, MIN, terminal, and cutoff cases. \hfill{\color{courseblue}\textbf{[2 points]}}
\item[\textbf{b.}] Trace a depth-two tree with at least two actions per internal node. \hfill{\color{courseblue}\textbf{[2 points]}}
\item[\textbf{c.}] Explain how depth is counted and how ties are broken. \hfill{\color{courseblue}\textbf{[1 points]}}
\end{enumerate}
\end{frame}

\begin{frame}{Task 4 - Alpha-beta and evaluation - part 1 (5 points)}
\small
Add pruning without changing minimax values and construct a defensible cutoff evaluation.

\vspace{0.2cm}
\begin{enumerate}
\item[\textbf{a.}] State the alpha and beta invariants and the pruning condition. \hfill{\color{courseblue}\textbf{[2 points]}}
\item[\textbf{b.}] Verify equality with plain minimax across generated small trees. \hfill{\color{courseblue}\textbf{[1 points]}}
\item[\textbf{c.}] Define an evaluation function as a weighted feature sum. \hfill{\color{courseblue}\textbf{[1 points]}}
\end{enumerate}
\end{frame}

\begin{frame}{Task 4 - Alpha-beta and evaluation - part 2 (5 points)}
\small
Add pruning without changing minimax values and construct a defensible cutoff evaluation.

\vspace{0.2cm}
\begin{enumerate}
\item[\textbf{d.}] Measure the effect of ordering and evaluation features. \hfill{\color{courseblue}\textbf{[1 points]}}
\end{enumerate}
\end{frame}

\begin{frame}{Task 5 - Integration and evidence (3 points)}
\small
Deliver an end-to-end agent and explain failures with reproducible evidence.

\vspace{0.2cm}
\begin{enumerate}
\item[\textbf{a.}] Run an end-to-end scenario that uses both pathfinding and adversarial choice. \hfill{\color{courseblue}\textbf{[1 points]}}
\item[\textbf{b.}] Test terminal roots, no legal actions, repeated states, and a narrow pruning window. \hfill{\color{courseblue}\textbf{[1 points]}}
\item[\textbf{c.}] Report nodes expanded, prune count, decision value, and runtime. \hfill{\color{courseblue}\textbf{[1 points]}}
\end{enumerate}
\end{frame}

\begin{frame}{Edge cases and defensive programming}
\small
\begin{itemize}
  \item The game is terminal at the root.
  \item A nonterminal state has no legal actions because of a domain defect.
  \item Different paths reach the same immutable state.
  \item All root actions have equal backed-up value.
  \item Alpha equals beta exactly after an update.
\end{itemize}
\end{frame}

\begin{frame}{Submission instructions}
\small
\begin{itemize}
  \item Submit the repository with modules, tests, and a reproducible demo command.
  \item Include a report with architecture, heuristic argument, evaluation features, and result tables.
  \item Provide a tagged final commit and exclude generated caches and credentials.
\end{itemize}
\vspace{0.2cm}
\alert{Do not submit credentials, generated caches, or unapproved libraries.}
\end{frame}

\begin{frame}{Grading rubric - part 1}
\scriptsize
\begin{tabularx}{\textwidth}{>{\bfseries}p{0.28\textwidth} p{0.08\textwidth} X}
\toprule
Category & Points & Required evidence \\
\midrule
Architecture and pathfinding & 4 & Architecture diagram, interfaces, and pathfinding tests. \\
A-star heuristic & 3 & Formula, proof sketch, and comparison against h=0. \\
Minimax decision engine & 5 & Minimax implementation, hand trace, and unit tests. \\
\bottomrule
\end{tabularx}
\end{frame}

\begin{frame}{Grading rubric - part 2}
\scriptsize
\begin{tabularx}{\textwidth}{>{\bfseries}p{0.28\textwidth} p{0.08\textwidth} X}
\toprule
Category & Points & Required evidence \\
\midrule
Alpha-beta and evaluation & 5 & Alpha-beta implementation, evaluation ablation, and pruning metrics. \\
Integration and evidence & 3 & Integrated demo, adversarial tests, and concise report. \\
\bottomrule
\end{tabularx}
\end{frame}

\begin{frame}{Reflection prompt}
In 200-250 words:

Explain one integration defect that looked like an algorithm defect. Describe the boundary contract, the diagnostic evidence, and the final repair.
\end{frame}

\end{document}
